\section{它是如何被证明的？}\label{sec:5}

在这一节，我描述进入 GLC 证明的一些主要想法。

\S\ref{sec:3.5} 的谱作用告诉我们，可以在 Hecke 算子的作用下对自守范畴 $\Shv(\Bun_{G})$ 作谱分解，可能的颜色（Langlands 参数）由 $\Gd$-局部系统、即特征簇 $\Loc_{\Gd}$ 的点提供。既然如此，我们就可以表述“最佳情形”：

\begin{itemize}
\item Hecke 算子有简单的联合谱（量子完全可积性），且
\item 每一种可能的颜色都出现，即谱恰好是整个 Langlands 参数空间 $\Loc_{\Gd}$。
\end{itemize}

换言之，我们可以期望所有联合本征空间的重数都为一，从而对每种颜色 $E\in\Loc_{\Gd}$ 都有唯一（在归一化意义下）的本征层 $\mathrm{Aut}_{E}\in\Shv(\Bun_{G})$。本征层条件是说，每个 Hecke 算子 $F\in\QCoh(\Loc_{\Gd})$ 通过乘以向量空间 $F_{E}$（即 $F$ 在 $E$ 处的纤维）作用在 $\mathrm{Aut}_{E}$ 上：
\[
F*\mathrm{Aut}_{E}\;\simeq\;F|_{E}\otimes\mathrm{Aut}_{E}.
\]
合起来，我们可以把这一期望表述为：谱作用 $\QCoh(\Loc_{\Gd})\curvearrowright\Shv(\Bun_{G})$ 是秩一的自由作用，于是归一化之后我们期望得到如下形式的谱定理：
\[
\Shv(\Bun_{G})\;\simeq\;\QCoh(\Loc_{\Gd}).
\]
回忆一下：线性算子有简单谱，当且仅当它有循环向量。因此我们可以把这一期望转述为：Hecke 作用有循环向量（确切说，循环对象），并且这个循环向量有满谱支撑，即它是所有颜色的叠加。换言之，我们要找一个对象 $\mathrm{White}_{G}\in\Shv(\Bun_{G})$，它具有白光的性质；一旦找到它，GLC 就应当由下面两条唯一刻画：

\begin{itemize}
\item \textbf{谱作用}：等价 $\Shv(\Bun_{G})\simeq\QCoh(\Loc_{\Gd})$ 保持 $\QCoh(\Loc_{\Gd})$ 的作用；
\item \textbf{归一化}：该等价互换 $\mathrm{White}_{G}$ 与 $\OO_{\Loc_{\Gd}}$。
\end{itemize}

\subsection{先做一点泛函分析}\label{sec:5.1}

这个“最佳情形”原来需要一些修正——也就是说，我们必须把泛函分析做对。

既然 GLC 关涉的是一种以层的范畴取代传统函数空间的调和分析，我们理应遇到熟悉的分析问题的回声——反过来说，GLC 为这种新的范畴化分析提供了试验场。与局部对称空间的调和分析一样，首要问题来自（拟）紧性的缺失——即，我们应当在尖点处施加什么样的增长条件？这些选择导向什么类的“拓扑向量空间”，它们带有何种对偶与算子理论？在范畴化背景下，这些问题又究竟意味着什么？

这些问题在 Gaitsgory 与 Drinfeld、Arinkin 合作的系列论文（包括~[DG13, DG15, Gai16, Gai17b, AG15]）中得到发展，其基础是 Lurie 发展的范畴机器~[Lur09, Lur]。一方面，他们确立了 de Rham GLC 中的基本函数空间，即范畴 $\Dmod(\Bun_{G})$，“奇迹般地”表现得像范畴化的希尔伯特空间——特别是它是自对偶的。另一方面，对谱侧的朴素猜测需要精化。在傅里叶理论中，紧支撑的缺失对偶于光滑性的缺失。因此，如果我们选择不对 $\Bun_{G}$ 上的层施加任何增长条件，就应当预期在 $\Loc_{\Gd}$ 上的层论中出现奇异性——即某种取值于向量空间的“分布”，而非“函数”。事实证明，层层面上的这一区别与底空间 $\Loc_{\Gd}$ 的奇异性密切相关。

这引出了代数几何中内容丰富的新篇章，即归纳凝聚层理论 $\IndCoh$~[GR17]。范畴 $\IndCoh$ 把 $\QCoh$ 扩大以容纳代数广义函数。与分布一样，这些新对象可以描述为在“测试”对象上的泛函——在这里是凝聚层上的函子。$\IndCoh$ 与 $\QCoh$ 的差别归结为：奇异簇上的凝聚层——例如奇异点处的摩天楼——不允许由向量丛给出的有限分解，或更几何地说，不能延拓到簇的任意无穷小形变。这引出了归纳凝聚层的奇异支撑这一概念——分布的波前集的精确对应物——它度量给定的广义层在哪些点与哪些形变方向上受到阻碍。

通过分析奇异支撑在层的自然操作下的行为，[AG15] 得以精确调校出：究竟何种类型的归纳凝聚层奇异性与自守侧给定的层论相匹配——我们应当允许幂零奇异支撑\anote{值得注意的是，这一概念与加在自守侧 Betti 层上的、名字相近的幂零支撑条件并无密切关系。}，这大致意味着不允许来自环面方向的任何奇异性（如同先前已理解的 $\mathrm{GL}_{1}$ 情形）。这最终导向非分歧几何 Langlands 猜想的精确陈述：它是一个等价
\[
\Shv(\Bun_{G})\;\simeq\;\IndCoh_{\Nilp}(\Loc_{\Gd})
\]
（de Rham 与 Betti 两种风味均如此），它把左侧 $\QCoh(\Loc_{\Gd})$ 的谱作用与谱侧自然的张量作用（在广义函数上乘以函数）等同起来。

\subsubsection{étale 背景}\label{sec:5.1.1}

里程碑式的论文~[AGK+20] 引入了 GLC 的一种新风味，即 étale（或受限）版本。它最重要的优点或许是：它在任意代数闭域上都有意义，从而第一次在正特征函数域的背景下表述了范畴化的 GLC，并建立了相应的谱作用。在 $\C$ 上，étale 版本构成 Betti 与 de Rham 猜想之间的中间“共同核心”。例如在谱侧，$\Loc^{\Betti}_{\Gd}$ 与 $\Loc^{\dR}_{\Gd}$ 有相同的点，且这些点有相同的形变理论，所以两者所有点的所有形式邻域之并是相同的。特征簇的 étale 版本恰好就是（对阿贝尔群而言）这个并，一般情形下则稍大一些（本质上是因为对幂零矩阵取指数是一个代数操作）。特别是，由摩天楼经有限过程搭出的凝聚层在三个版本上都是相同的。对偶地，étale 自守范畴被设计成由单色对象离散地、而非连续地搭成。

论文~[AGK+20] 证明了 GLC 的 Betti 与 de Rham 版本彼此等价：它表明两者都化归到共同的 étale 版本。这让我们可以交替利用两种背景的优势——例如，Betti 特征栈 $\Loc^{\Betti}_{\Gd}$ 非常简单，是仿射簇对 $\Gd$ 的商，明显地只依赖于 $C$ 的拓扑；而 de Rham 版本 $\Loc^{\dR}_{\Gd}$ 远非仿射，于是凝聚层之间的 Hom 空间是有限维的。

\subsection{间奏：相对 Langlands，或 $L$-函数都去哪儿了？}\label{sec:5.2}

Langlands 哲学所预言的并不只是自守侧与谱侧之间使 Hecke 算子对角化的一个关系：它还配备了庞大的预测体系，说明什么样的测量应当与什么样的测量相匹配。对这些测量的研究称为相对 Langlands 纲领，并在~[SV17] 中被系统化。最著名的一项——把自守表示与 Galois 表示相联系的 $L$-函数的匹配——可以说是算术 Langlands 纲领的主要特征。$L$-函数本身在 GLC 的大部分文献中显著地缺席，但类似的匹配测量在归一化、刻画与证明 GLC 中起着主要作用。

自守形式的“测量”是什么意思？它们是自守形式向量空间上的线性泛函——不过我们也可以（并且将会）对偶地把它看作与某个特定向量的内积，或在几何语境中看作从某个特定层出发的 Hom 空间。我们可以为之作出预测的特殊测量——周期（period）\anote{容易混淆的是，不要与代数几何中的“周期”混为一谈。}——是高度结构化的：它们在 Langlands 纲领的所有不同背景（局部、整体、算术、几何）中都有对应物，并满足很强的相容性。

在物理中，这种结构化的测量或状态集合是任何量子场论的本质部分，它们对应于边界理论（或余维一的缺陷）。更具启发性的是，我们可以通过规范场论与带电物质的相互作用来研究它。在电磁 S-对偶的语境中，Gaiotto–Witten~[GW09b, GW09a] 开创了这类研究，并在~[BZSV24] 中被提议为相对 Langlands 纲领的一般框架（另见~[Gan]）。\zhnote{“周期”在 Langlands 语境中不应与代数几何的周期（周期积分）混淆：它指自守形式上的某种不变线性泛函，即把自守表示“求值／积分”得到一个数或向量空间。相对 Langlands 纲领（Relative Langlands）就是系统研究这些周期、并把它们与谱侧相应的“测量”相匹配的理论；nLab 词条 \emph{relative Langlands program} 给出的正是这一观点：\href{https://ncatlab.org/nlab/show/relative+Langlands+program}{\texttt{https:/\allowbreak{}/\allowbreak{}ncatlab.\allowbreak{}org/\allowbreak{}nlab/\allowbreak{}show/\allowbreak{}relative+\allowbreak{}Langlands+\allowbreak{}program}}（最后修订 2025-04-29）。}

一个基本例子是平凡周期——常数波或零频波，它对应于最简单的测量：取自守形式的平均，或取自守层的整体上同调。这个周期在 GLC 的语境中已经非常有趣：它在~[GL19b] 中得到研究，并被用来证明函数域的 Weil 玉河数猜想。

下面三节中，我们来看三种重要的测量（周期），它们的相互作用构成了 GLC 证明的骨架（转述自~[CCF+24]）。这里列出相匹配的自守测量与谱测量的名称及其主要作用：

\begin{itemize}
\item \textbf{Eisenstein $\leftrightarrow$ Eisenstein}：在 GLC 上提供归纳结构，把猜想化归到尖点情形。
\item \textbf{Whittaker $\leftrightarrow$ 整体截面（Neumann）}：用于归一化——提供对尖点自守层的忠实测量。
\item \textbf{Kac–Moody $\leftrightarrow$ opers}：用于构造尖点自守层。
\end{itemize}

周期的可能最关键的特征是：它们满足很强的局部–整体相容性——它们由 $\Theta$-级数给出。$\Theta$-级数（或 Poincaré 级数）是经典 Riemann $\theta$-函数的推广，后者例如出现在 Riemann 关于 $\zeta(s)$ 的积分公式中。它们是模形式与自守形式，写成局部数据在一个整体对称群上的平均。这事实上是写出自守形式最常见的方式之一。

$L$-函数的一个特征性质是它们作为 Euler 乘积的“局部–整体”描述。在几何背景中，周期有一个松散类比的关键性质——又是因子化。也就是说，我们可以把 $\Theta$-级数附着到曲线 $C$ 上有限多个点处的局部数据上，而构造与点的碰撞相容。这一因子化性质是把局部输入装配成可理解的整体输出的胶水。它在技术上也极具挑战——GLC 证明（及其许多前身）之所以冗长，一个主要原因是需要整体地理解所有构造的局部、整体与可因子化版本。

\subsection{归纳步：Eisenstein 级数}\label{sec:5.3}

Eisenstein 级数同时是自守形式理论中困难的部分和容易的部分。一方面，对分析的严肃需求来自模曲线（以及更一般的算术局部对称空间）的连续谱，这是它们非紧性的效应——即来自在无穷远处（即在尖点处）不消失的函数。Eisenstein 级数的一般解析理论是 Langlands 的得意成就之一，它给出了这一连续谱的构造。Eisenstein 级数的几何 Langlands 对应物必须在范畴化分析的勇敢新世界中，应对与经典理论中的解析困难平行的问题。

另一方面，从表示论看，Eisenstein 级数提供了理论中“不出人意料”的部分——由抛物归纳而来的部分。这是最高权理论的精致化，让我们能从较小 Levi 子群（对 $\mathrm{GL}_{n}$ 即对角或分块对角矩阵构成的子群）的表示搭出大群的表示。约化群表示论的指路明灯——Harish-Chandra 的尖点形式哲学——说：对每个群，我们应首先筛出由较小群经抛物归纳而来的表示族（或系列），只留下稀有的珍宝——尖点表示——然后带着它们走向越来越大的群。

整个 Langlands 纲领都遵循这一归纳策略：它规定，从 Levi 子群 $M\subset G$ 构造 Eisenstein 级数，应当与对偶 Levi 子群 $\check M\subset\Gd$ 的类似抛物归纳相匹配，这恰好解释了可约 Langlands 参数所在的轨迹。因此粗略地说，我们需要先“减去”对应两侧 Eisenstein 级数的贡献，只留下尖点自守函数（那些 Eisenstein 周期消失的函数），然后把它们与不可约 Langlands 参数相匹配。

把漫长而技术性的故事压缩到极短：沿着这些一般模式，经过 Drinfeld、Laumon、Braverman、Gaitsgory、Finkelberg、Mirkovic、Bezrukavnikov 等人数十年的贡献，Eisenstein 级数的理论如今在几何背景中已得到很好的理解。专著~[CCF+24] 完成了归纳步，特别是证明了 Eisenstein 级数与如上 Langlands 函子 $L$ 的相容性。有了这一结果，GLC 的证明就化归为对 Eisenstein 级数正交补——尖点自守层——的研究。

\subsection{白光：Whittaker 周期}\label{sec:5.4}

与单色波相对的另一极端是白光，它是一切颜色的叠加。在傅里叶变换的框架中，白光由 $x=0$ 处的 $\delta$-函数 $\delta_{0}$ 表示，其傅里叶变换是常函数 $1$。这一选择在归一化微分（或平移）的本征向量时起本质作用，因为本征性质只在相差一个标量的意义下刻画它们。归一化的指数函数被选为在 $x=0$ 处取值 $1$，即通过与 $\delta$-函数 $\delta_{0}$ 的配对来选取，从而成为白光的组分。

这一归一化很容易推广到更一般的阿贝尔对偶——我们用在本元处的取值来归一化本征函数。换言之，我们规定对偶把本元处的 $\delta$-函数 $\delta_{e}$（一个 Dirichlet 边界条件）送到常函数 $1$（一个 Neumann 边界条件）。图~\ref{fig:1}（梅林变换情形）中对应白光的那个白色极点展示了这一点。

在非阿贝尔对偶中，表示论的经验提示，白光的角色由一个不那么熟悉的客体扮演，即 Whittaker 周期\anote{以数学家 Sir Edmund Whittaker 命名。机缘凑巧的是，Whittaker 这个姓氏在古英语中的词源意为“White Field”（白色田野）。}（在物理中称为 Nahm 极点边界条件~[GW09a]）。换言之，我们希望这样来归一化 Langlands 对应：规定 Whittaker 周期被送到由整体积分或整体截面给出的谱测量（即送到 Neumann 边界条件）。也就是说，所求的对偶应当把表示对象——Whittaker 层——送到表示谱侧整体截面的结构层 $\OO_{\Loc_{\Gd}}$。

这个 Whittaker 周期是什么？它来自 Whittaker 函数——数学物理中 Whittaker 微分方程的解——这些函数曾被用来给出局部域上群的表示的归一化描述。在模形式的框架中，Whittaker 周期对应于 $q$-展开中第一个傅里叶系数的测量，而后者正是用来归一化尖点形式的。在 GLC 的框架中，Whittaker 周期有一个引人入胜的微局部几何描述（见图~\ref{fig:3}）——它由 Hitchin 截面~[Hit92] 表示。这是 Hitchin 可积系统的一个截面，由那些作为矩阵处于有理标准形（对 $G=\mathrm{GL}_{n}$，或更一般地，位于 Kostant 对有理标准形所作的约化群推广中）的希格斯场构成。它们对应于配备平凡线丛的谱曲线，即 Hitchin 纤维的单位元。与阿贝尔情形的一个关键差别（图中也可见）是：Hitchin 截面完全错过零截面 $\Bun_{G}\subset T^{*}\Bun_{G}$（因为零矩阵不容许循环向量）。

注意，作为截面，它与 Hitchin 纤维横截——而这正是我们关于本征层长什么样的微局部图景所期望的，也正是一个白光候选者所应有者。

一旦挑出 Whittaker 层，我们就可以开始用 Hecke 函子作用在它上面。（微局部地，这意味着沿 Hitchin 纤维平移来移动 Hitchin 截面。）借助谱作用，这一指派确定一个函子
\[
\QCoh(\Loc_{\Gd})\longrightarrow\Shv(\Bun_{G}),\qquad F\longmapsto F*\mathrm{Whit}_{G},
\]
它把结构层 $\OO_{\Loc_{\Gd}}$ 送到 Whittaker 层。如果我们真的期望自守侧在谱作用下是秩一的自由模，那么这就是 GLC 了。

但我们知道这个函子不可能完全正确：我们需要把谱侧修正为稍大的范畴 $\IndCoh_{\Nilp}(\Loc_{\Gd})$。这需要一些细致的分析。[FrR25, NT25] 利用微局部解释，证明了 Whittaker 系数（由 Whittaker 层表示的测量）一些出人意料的好性质，包括在一个大的层子范畴上的正合性。这些结果与~[GR24a] 中进一步的估计（关于上同调有界性与范畴紧性）相结合，表明：用形如 $F*\mathrm{Whit}_{G}$（$F\in\QCoh(\Loc_{\Gd})$）的层来测试自守层，可以唯一地提升为反方向的函子
\[
\Shv(\Bun_{G})\longrightarrow\IndCoh_{\Nilp}(\Loc_{\Gd}).
\]
Whittaker 周期的微局部几何还被~[FrR25] 用来证明几何 Langlands 对应的一个关键单射性质，它与数论中的经验多少相反：Whittaker 系数在尖点层上不消失。也就是说，我们可以用 Whittaker 层及其 Hecke 平移忠实地测试任何尖点层，所以它不会被函子 $L$ 湮灭。结合对 Eisenstein 部分的归纳理解，剩下的就是设法建立满射性——看到每种颜色都出现。

\subsection{构造：Kac–Moody 局部化与 opers}\label{sec:5.5}

最后，我们需要找到一些稀有的珍宝——尖点自守层。事实上需要很多，以纳入所有不可约 Langlands 参数。如我们提到的，实际写出经典自守形式的主要技术是 $\Theta$-级数理论——即从局部数据出发再作平均。Beilinson 与 Drinfeld 发现的 de Rham 几何 Langlands 的王牌，是一类特别的几何 $\Theta$-级数，即 Kac–Moody 周期，它有一个以 opers 给出的具体对偶描述。这一构造与我们讨论过的所有话题都不同：它在 Langlands 纲领的其他部分没有任何已知的直接对应物。

它为何如此难以捉摸？答案是：整个 de Rham GLC 建立在李代数的表示论而非李群的表示论之上，而 Kac–Moody 周期尤其关涉不能积分到群上的表示，例如 $\mathfrak{g}$ 在其包络代数 $U\mathfrak{g}$ 自身上的作用。这类表示在 Betti 或 étale GLC 中没有对应物，在 $p$-进群的平滑表示论（局部 Langlands 对应的主题）中也没有\anote{不过在 $p$-进群的局部解析（locally analytic）表示这一迅速发展的领域中有对应物。}。

另一种形式的答案来自物理：Kac–Moody 周期不是拓扑的——它对黎曼面 $C$ 的复结构敏感。因此它关联于远为丰富的共形（或全纯）量子场论的研究，而不是相对已被很好理解的拓扑场论~[FG20]。

那么这个 Kac–Moody 周期讲的是什么？整体版本是微分算子层 $\mathcal{D}_{\Bun_{G}}$ 本身——一个非常自然的对象，但它在单值性与可构造层的拓扑世界中没有任何类似物。更确切地说，我们需要用涉及自旋结构的微分算子修正，这是量子化的普遍特征。局部版本是环路李代数 $\mathfrak{g}((z))$——即 $\LG$ 的李代数——的表示范畴，同样差一个“自旋”修正，给出所谓的 $\mathfrak{g}((z))$ 的中心扩张（仿射 Kac–Moody 代数 $\widehat{\mathfrak{g}}$）的临界水平表示。局部–整体构造（$\Theta$-级数）由几何表示论最伟大的成就之一给出，即 Beilinson–Bernstein 局部化构造：一种由李代数的表示写出微分方程组的方法。通过这一构造，表示论提供了 $\Bun_{G}$ 上大量相当具体的 D-模；Beilinson 与 Drinfeld 意识到，这正是建立 Langlands 对应的秘方。

表示论在这里能提供什么？这个配方的起手式是 Feigin 与 Frenkel 的一个奇妙定理~[FF92]（另见专著~[Fre07]），它使用了共形场论即顶点代数理论所提示的环路代数途径。该定理把 $\widehat{\mathfrak{g}}$ 的临界水平表示的内蕴参数，等同为去心圆盘上一类特殊的 Langlands 参数（$\Gd$-局部系统），即 opers。opers（在 $\mathrm{GL}_{n}$ 情形）是选定了微分循环向量的联络——即处于有理标准形的联络。它们恰好就是当我们想解单个 $n$ 阶常微分方程时得到的一个变量中 $n$ 个一阶微分方程组。opers 构成 Whittaker 周期的对偶、谱的对应物——是平坦联络空间 $\Loc_{\Gd}$ 上 Hitchin 截面的一个版本。

Feigin–Frenkel 定理提供了某种非凡的东西——一座直接沟通与 $G$ 相关的表示论和 Langlands 对偶侧的桥梁，我们可以用它往自守侧填充本征层。如 \S\ref{sec:4.2} 所述，Beilinson–Drinfeld 写出了 $\Bun_{G}$ 上一族由整体微分算子给出的显式微分方程——即直接由 $\mathcal{D}_{\Bun_{G}}$ 搭成——它们是 \S\ref{sec:1.3} 中 $C^{\times}$ 上 $z\frac{d}{dz}f=sf$ 及其在雅可比簇上的对应物（\S\ref{sec:2.4.1}）的非阿贝尔类比。通过把这些 D-模辨识为从对应于整条曲线 $C$ 上 $\Gd$-opers 的那些 $\widehat{\mathfrak{g}}$-表示出发所作局部–整体构造的输出，他们得以证明：这些 D-模确实是 Hecke 本征层，其本征值恰是这些 opers（看作 $\Loc_{\Gd}$ 的点）。

但这只是开始：opers 本身只构成 $\Loc_{\Gd}$ 的一个半维子簇。然而任何不可约局部系统都在有限多个点之外容许 oper 结构~[Ari16]（对 $\mathrm{GL}_{n}$ 这是 Deligne 的经典结果：平坦联络容许亚纯循环向量）。而且同伦意义下这样的选择只有一个——选择空间是可缩的（对经典群是~[BKS22] 的定理，一般情形则是 GLC 的推论）。于是我们必须向表示论要求更多：对带有有理 oper 结构的 $\Gd$-局部系统所关联的一类更广的表示给出描述，这由~[FG09a] 提供。现在可以下锅了。

在~[ABC+24b] 中，临界水平表示论被搭成一台高度结构化的机器，把因子化与局部–整体相容性纳入其中。这台机器被用来往自守侧填充来自 Kac–Moody 局部化的 D-模，往谱侧填充来自带 oper 结构的局部系统的相匹配的层。这给出了足够丰富的匹配对象备货，以看出 Langlands 函子的“满射性”。

把这些配料合在一起，我们就有了五篇论文组成的序列：

\begin{itemize}
\item [GR24a]：Langlands 函子的构造；
\item [ABC+24b]：用 Kac–Moody 局部化构造尖点层；
\item [CCF+24]：由 Eisenstein 级数给出的归纳步；
\item [ABC+24a]：对 Langlands 函子更细致的泛函分析（它满足范畴自对偶性，即 ambidexterity）；
\item [GR24b]：得出所有本征空间重数为一、非分歧 GLC 成立的结论。
\end{itemize}
